%%%%%%%%%%%%%%%%%%%%%%%%%%%%%%%%%%%%%%%%%%%%%%%%%%%%%%%%%%%%%%%%%%%%%%
% Lớp 8 | Chương X. Một số hình khối trong thực tiễn | Hình chóp tứ giác đều
% Dạng 8H0?39-3: Bài toán thực tế
%%%%%%%%%%%%%%%%%%%%%%%%%%%%%%%%%%%%%%%%%%%%%%%%%%%%%%%%%%%%%%%%%%%%%%


% nguồn: De_AI_Toan8_cuoi_hoc_ki_i_1 (nhập 21/11/2026)
\begin{bt}%[8H0H39-3]%[ma:5c678080]
	Một chiếc lều cắm trại có dạng hình chóp tứ giác đều với mặt đáy là hình vuông cạnh $3\text{ m}$ và chiều cao của lều là $2\text{ m}$.
a) Tính thể tích không khí bên trong chiếc lều.
b) Tính diện tích vải bạt cần dùng để làm bốn mặt bên của lều (không tính phần vải thừa ở mép may).
	\loigiai{
		a) Diện tích đáy của lều là $S_{\text{đáy}} = 3^2 = 9\text{ m}^2$. \\
Thể tích không khí bên trong lều là $V = \dfrac{1}{3} \cdot 9 \cdot 2 = 6\text{ m}^3$. \\
b) Khoảng cách từ tâm đáy đến cạnh đáy là $r = \dfrac{3}{2} = 1{,}5\text{ m}$. \\
Độ dài trung đoạn của lều là $d = \sqrt{2^2 + 1{,}5^2} = \sqrt{4 + 2{,}25} = \sqrt{6{,}25} = 2{,}5\text{ m}$. \\
Diện tích vải bạt làm bốn mặt bên là $S_{xq} = (2 \cdot 3) \cdot 2{,}5 = 15\text{ m}^2$.
	}
\end{bt}

% nguồn: De_AI_Toan8_cuoi_hoc_ki_i_1 (nhập 21/11/2026)
\begin{bt}%[8H0H39-3]%[ma:b7753d1b]
	Một công trình kiến trúc mô phỏng kim tự tháp bằng kính có dạng hình chóp tứ giác đều với cạnh đáy dài $30\text{ m}$ và chiều cao $20\text{ m}$.
a) Tính thể tích không gian bên trong kim tự tháp.
b) Tính tổng diện tích kính cần dùng để lắp kín bốn mặt bên của kim tự tháp (bỏ qua diện tích các khung thép).
	\loigiai{
		a) Diện tích mặt đáy là $S_{\text{đáy}} = 30^2 = 900\text{ m}^2$. \\
Thể tích không gian bên trong kim tự tháp là $V = \dfrac{1}{3} \cdot 900 \cdot 20 = 6000\text{ m}^3$. \\
b) Khoảng cách từ tâm đáy đến cạnh đáy là $r = \dfrac{30}{2} = 15\text{ m}$. \\
Độ dài trung đoạn là $d = \sqrt{20^2 + 15^2} = \sqrt{400 + 225} = \sqrt{625} = 25\text{ m}$. \\
Diện tích kính cần dùng là $S_{xq} = (2 \cdot 30) \cdot 25 = 1500\text{ m}^2$.
	}
\end{bt}

% nguồn: De_AI_Toan8_cuoi_hoc_ki_i_1 (nhập 21/11/2026)
\begin{bt}%[8H0H39-3]%[ma:08a76794]
	Phần mái của một chòi nghỉ mát trong công viên có dạng hình chóp tứ giác đều với cạnh đáy dài $4\text{ m}$ và trung đoạn dài $2{,}5\text{ m}$.
a) Tính chiều cao của phần mái chòi.
b) Tính diện tích phần mái cần lợp ngói và chi phí lợp mái, biết đơn giá thi công là $350\,000\text{ đồng/m}^2$.
	\loigiai{
		a) Khoảng cách từ tâm đáy đến cạnh đáy là $r = \dfrac{4}{2} = 2\text{ m}$. \\
Chiều cao phần mái chòi là $h = \sqrt{2{,}5^2 - 2^2} = \sqrt{6{,}25 - 4} = \sqrt{2{,}25} = 1{,}5\text{ m}$. \\
b) Diện tích phần mái cần lợp ngói chính là diện tích xung quanh của hình chóp: \\
$S_{xq} = (2 \cdot 4) \cdot 2{,}5 = 20\text{ m}^2$. \\
Tổng chi phí lợp mái là $20 \cdot 350\,000 = 7\,000\,000\text{ đồng}$.
	}
\end{bt}

% nguồn: De_AI_Toan8_cuoi_hoc_ki_i_1 (nhập 21/11/2026)
\begin{bt}%[8H0H39-3]%[ma:6c01b241]
	Một cửa hàng đặt làm $50$ hộp quà có dạng hình chóp tứ giác đều (có cả nắp đáy) với cạnh đáy $12\text{ cm}$ và chiều cao $8\text{ cm}$.
a) Tính diện tích toàn phần của một hộp quà.
b) Biết phần bìa cứng bị hao hụt khi cắt gấp chiếm $10\%$ diện tích toàn phần của các hộp, tính tổng diện tích bìa cứng cần dùng để làm $50$ hộp quà.
	\loigiai{
		a) Khoảng cách từ tâm đáy đến cạnh đáy là $r = \dfrac{12}{2} = 6\text{ cm}$. \\
Độ dài trung đoạn của hộp quà là $d = \sqrt{8^2 + 6^2} = 10\text{ cm}$. \\
Diện tích xung quanh một hộp quà: $S_{xq} = (2 \cdot 12) \cdot 10 = 240\text{ cm}^2$. \\
Diện tích đáy: $S_{\text{đáy}} = 12^2 = 144\text{ cm}^2$. Diện tích toàn phần một hộp quà: $S_{tp} = 240 + 144 = 384\text{ cm}^2$. \\
b) Diện tích toàn phần của $50$ hộp quà là $50 \cdot 384 = 19\,200\text{ cm}^2$. \\
Tính thêm $10\%$ hao hụt, tổng diện tích bìa cứng cần dùng là $19\,200 \cdot 1{,}1 = 21\,120\text{ cm}^2$.
	}
\end{bt}

% nguồn: De_AI_Toan8_cuoi_hoc_ki_i_1 (nhập 21/11/2026)
\begin{bt}%[8H0H39-3]%[ma:db229c35]
	Mỗi chiếc bánh ít lá gai truyền thống có dạng hình chóp tứ giác đều với cạnh đáy bằng $6\text{ cm}$ và chiều cao bằng $5\text{ cm}$.
a) Tính thể tích của một chiếc bánh.
b) Nếu làm một mẻ gồm $25$ chiếc bánh như nhau thì tổng thể tích bánh là bao nhiêu xăng-ti-mét khối?
	\loigiai{
		a) Diện tích mặt đáy của chiếc bánh là $S_{\text{đáy}} = 6^2 = 36\text{ cm}^2$. \\
Thể tích một chiếc bánh là $V = \dfrac{1}{3} \cdot 36 \cdot 5 = 60\text{ cm}^3$. \\
b) Tổng thể tích của $25$ chiếc bánh là $25 \cdot 60 = 1500\text{ cm}^3$.
	}
\end{bt}

% nguồn: De_AI_Toan8_cuoi_hoc_ki_i_1 (nhập 21/11/2026)
\begin{bt}%[8H0H39-3]%[ma:ccc092fd]
	Một cột mốc bằng bê tông gồm hai phần: phần dưới là một hình lập phương có cạnh $6\text{ dm}$, phần trên là một hình chóp tứ giác đều có mặt đáy trùng với mặt trên của hình lập phương và chiều cao bằng $4\text{ dm}$.
a) Tính thể tích của cả cột mốc bê tông.
b) Người ta sơn toàn bộ mặt ngoài của cột mốc (trừ mặt đáy dưới cùng tiếp xúc với mặt đất). Tính diện tích bề mặt cần sơn.
	\loigiai{
		a) Thể tích phần hình lập phương là $V_1 = 6^3 = 216\text{ dm}^3$. \\
Thể tích phần hình chóp tứ giác đều là $V_2 = \dfrac{1}{3} \cdot 6^2 \cdot 4 = 48\text{ dm}^3$. \\
Thể tích cả cột mốc là $V = V_1 + V_2 = 216 + 48 = 264\text{ dm}^3$. \\
b) Diện tích $4$ mặt bên của hình lập phương là $S_1 = 4 \cdot 6^2 = 144\text{ dm}^2$. \\
Đối với phần hình chóp tứ giác đều, độ dài trung đoạn là $d = \sqrt{4^2 + 3^2} = 5\text{ dm}$. \\
Diện tích xung quanh phần hình chóp là $S_2 = (2 \cdot 6) \cdot 5 = 60\text{ dm}^2$. \\
Tổng diện tích bề mặt cần sơn là $S = S_1 + S_2 = 144 + 60 = 204\text{ dm}^2$.
	}
\end{bt}
